% TEMPLATE-MILC-v3.tex - plantilla maestra para todas las guias MILC v3
% Ver scripts/generadores/build-guia-milc.py para la lista completa de
% claves esperadas. El script valida que no queden placeholders sin
% reemplazar antes de escribir el archivo final.

\documentclass[11pt,letterpaper]{article}
\usepackage[margin=1.75cm]{geometry}
\usepackage[table]{xcolor}
\usepackage{fontspec}
\usepackage[spanish]{babel}
\usepackage{tikz}
\usepackage[most]{tcolorbox}
\usepackage{tabularx}
\usepackage{array}
\usepackage{fancyhdr}
\usepackage{titlesec}
\usepackage{ragged2e}
\usepackage{enumitem}
\usepackage{microtype}

\setmainfont{Helvetica Neue}
\setsansfont{Helvetica Neue}
\setlength{\parindent}{0pt}
\setlength{\parskip}{6pt}
\hyphenpenalty=200
\exhyphenpenalty=200
\pretolerance=400
\tolerance=2000
\setlength{\emergencystretch}{1em}
\renewcommand{\arraystretch}{1.25}
\setlist{leftmargin=5mm,itemsep=1mm,topsep=1mm}
\newcolumntype{Y}{>{\RaggedRight\arraybackslash}X}

\definecolor{milcMagenta}{HTML}{E6007E}
\definecolor{milcVino}{HTML}{5A0038}
\definecolor{milcTurquesa}{HTML}{0093A5}
\definecolor{milcVerde}{HTML}{58A923}
\definecolor{milcMostaza}{HTML}{E5B400}
\definecolor{milcOcre}{HTML}{B86600}
\definecolor{milcNegro}{HTML}{241816}
\definecolor{milcGris}{HTML}{F7F7F4}
\definecolor{milcLinea}{HTML}{E8E2DD}
\definecolor{milcGradePrimary}{HTML}{0066FF}
\definecolor{milcGradeDark}{HTML}{003D99}
\definecolor{milcGradeSoft}{HTML}{D6E8FF}
\definecolor{milcGradeAccent}{HTML}{CFE7FF}
\definecolor{milcGradeText}{HTML}{FFFFFF}
\color{milcNegro}

\pagestyle{fancy}
\fancyhf{}
\lhead{\small Tecnología e Informática · Grado 11}
\rhead{\small Guía 19 · MILC v3}
\cfoot{\small\thepage}
\renewcommand{\headrulewidth}{0pt}

\newtcolorbox{titlebox}[1]{
  enhanced,colback=#1,colframe=#1,arc=7mm,boxrule=0pt,
  left=7mm,right=7mm,top=3mm,bottom=3mm,
  before skip=8pt,after skip=8pt,
  fontupper=\sffamily\bfseries\Large\color{white}
}

\newtcolorbox{softbox}[3]{
  enhanced,colback=#2,colframe=#1,borderline west={4pt}{0pt}{#1},
  arc=2mm,boxrule=.4pt,left=4mm,right=4mm,top=3mm,bottom=3mm,
  before skip=6pt,after skip=8pt,title={#3},coltitle=white,
  colbacktitle=#1,fonttitle=\sffamily\bfseries,fontupper=\RaggedRight,
  attach boxed title to top left={xshift=3mm,yshift=-2mm},
  boxed title style={arc=2mm, boxrule=0pt}
}

\newtcolorbox{drawbox}[2]{
  enhanced,colback=white,colframe=#1,arc=4mm,boxrule=1pt,
  left=5mm,right=5mm,top=4mm,bottom=4mm,before skip=8pt,after skip=8pt,
  title={#2},coltitle=white,colbacktitle=#1,fonttitle=\sffamily\bfseries,
  fontupper=\RaggedRight,
  attach boxed title to top left={xshift=4mm,yshift=-2mm},
  boxed title style={arc=3mm, boxrule=0pt}
}

\newtcolorbox{infoband}[1]{
  enhanced,colback=#1!8,colframe=#1!50,arc=2mm,boxrule=.5pt,
  left=4mm,right=4mm,top=2mm,bottom=2mm,before skip=4pt,after skip=4pt,
  fontupper=\small\RaggedRight
}

% Puente narrativo entre secciones (contrato editorial Conéctate).
% Da continuidad: "Acabas de X. Ahora vas a Y porque Z."
% Visualmente: banda izquierda en color del grado, texto en itálica pequeña.
\newtcolorbox{puentebox}{
  enhanced,colback=milcGradeSoft,colframe=milcGradeDark,
  borderline west={3pt}{0pt}{milcGradeDark},
  arc=0mm,boxrule=0pt,left=5mm,right=4mm,top=2.5mm,bottom=2.5mm,
  before skip=4pt,after skip=8pt,
  fontupper=\itshape\small\color{milcNegro}\RaggedRight
}
\newcommand{\puente}[1]{%
  \begin{puentebox}%
    \textcolor{milcGradeDark}{\textbf{\textrightarrow}}\hspace{2mm}#1%
  \end{puentebox}%
}

\newcommand{\linead}[1]{\noindent\color{milcLinea}\rule{#1}{0.5pt}\color{milcNegro}}
\newcommand{\respuesta}[1]{\par\vspace{2mm}\foreach \n in {1,...,#1}{\noindent\color{milcLinea}\rule{\linewidth}{0.5pt}\par\vspace{5mm}}\color{milcNegro}}
\newcommand{\checkbox}{\textcolor{milcGradeDark}{\fbox{\phantom{x}}}\hspace{5pt}}

\newcommand{\phasecard}[4]{
\begin{tcolorbox}[enhanced,colback=#3,colframe=#2,arc=3mm,boxrule=.5pt,height=3.05cm,valign=center,halign=center,left=1.5mm,right=1.5mm,top=2mm,bottom=2mm]
{\sffamily\bfseries\fontsize{22}{24}\selectfont\textcolor{#2}{#1}}\par
{\sffamily\bfseries\small #4}
\end{tcolorbox}
}

\begin{document}
\thispagestyle{empty}
\begin{tikzpicture}[remember picture,overlay]
  \fill[white] (current page.south west) rectangle (current page.north east);
  \path[fill=milcGradePrimary,rounded corners=16mm]
    ([xshift=.55cm,yshift=-.55cm]current page.north west) rectangle
    ([xshift=-.55cm,yshift=3.65cm]current page.south east);

  \node[anchor=north west,text=milcGradeText] at ([xshift=2.0cm,yshift=-1.85cm]current page.north west)
    {\sffamily\bfseries\fontsize{14}{16}\selectfont GRADO};
  \node[anchor=north west,text=milcGradeText,opacity=.82] at ([xshift=2.0cm,yshift=-2.75cm]current page.north west)
    {\sffamily\bfseries\fontsize{9}{11}\selectfont SERIE GUÍAS · TECNOLOGÍA E INFORMÁTICA};

  \begin{scope}[shift={([xshift=-3.05cm,yshift=-2.55cm]current page.north east)}]
    \fill[white,opacity=.80] (-.62,-.52) rectangle (.62,.52);
    \draw[milcGradeText,line width=1pt] (-.62,-.52) rectangle (.62,.52);
    \fill[milcGradeDark] (-.42,-.38) rectangle (-.22,.06);
    \fill[milcGradeAccent] (-.10,-.38) rectangle (.10,.24);
    \fill[milcGradePrimary!60!black] (.22,-.38) rectangle (.42,.38);
  \end{scope}

  \node[anchor=west,text=milcGradeText] at ([xshift=1.9cm,yshift=-12.85cm]current page.north west)
    {\sffamily\bfseries\fontsize{124}{124}\selectfont 11};
  \draw[milcGradeText,line width=1.7pt]
    ([xshift=4.52cm,yshift=-11.68cm]current page.north west) circle (.13cm);

  \node[anchor=west,text=milcGradeText,text width=8.7cm,align=left] at ([xshift=10.35cm,yshift=-11.6cm]current page.north west)
    {
     {\sffamily\bfseries\fontsize{19}{22}\selectfont Mejora continua ---\\Kaizen aplicado a\\procesos digitales}\\[4mm]
     {\sffamily\bfseries\fontsize{23}{26}\selectfont Guía 19-11-TIC}\\[2mm]
     {\sffamily\fontsize{13}{16}\selectfont Período 2 · Automatización y procesos}\\[5mm]
     {\sffamily\bfseries\fontsize{11}{13}\selectfont Institución Educativa Sor María Juliana}\\[1mm]
     {\sffamily\fontsize{10}{12}\selectfont Autor: Dr. Álvaro Cárdenas Orozco}
    };

  \path[fill=milcGradeSoft,rounded corners=7mm]
    ([xshift=1.35cm,yshift=.85cm]current page.south west) rectangle
    ([xshift=-1.35cm,yshift=4.25cm]current page.south east);
  \draw[milcGradeDark,line width=.65pt,rounded corners=7mm]
    ([xshift=1.35cm,yshift=.85cm]current page.south west) rectangle
    ([xshift=-1.35cm,yshift=4.25cm]current page.south east);
  \node[anchor=center,text=milcNegro,text width=17.0cm,align=center] at ([yshift=2.55cm]current page.south)
    {\sffamily\fontsize{10}{12}\selectfont Metodología MILC · Apertura ancestral · Pensamiento computacional · Triángulo Dussel-Estoicismo-Floridi\\
    \textcolor{milcGradeDark}{\bfseries Producto:} 1 propuesta de 3 mejoras concretas al proceso del periodo, cada mejora con descripción operativa, impacto esperado (alto/medio/bajo), esfuerzo requerido (alto/medio/bajo), priorización en matriz impacto-esfuerzo y plan de implementación en 2 frases};
\end{tikzpicture}
\mbox{}
\newpage

\section*{Apertura: Mejora continua --- Kaizen aplicado a procesos digitales}

\begin{softbox}{milcGradeDark}{milcGradeSoft}{Saber ancestral}
El oficio se afina con repetición y crítica: el \textbf{maestro carpintero} de Cartago miraba la silla terminada y veía qué podría haber salido mejor; el siguiente día corregía el ajuste de la pata o el grano del corte. La \textbf{tejedora Wayuu} aprendía con cada manta a apretar más uniforme, a elegir mejor el hilo, a calcular los rombos con más precisión. La \textbf{cocinera mayor} ajustaba el sancocho receta tras receta: \emph{``hoy le falta sal''}, \emph{``hoy se pegó el grano''}. Nunca se daba por terminado el oficio: cada pieza era un poco mejor que la anterior. Los japoneses le pusieron nombre técnico a esta sabiduría: \textbf{Kaizen} (改善, \emph{cambio hacia mejor}), pero el gesto existía mucho antes en el saber del oficio.
\end{softbox}

\begin{softbox}{milcTurquesa}{milcGris}{Saber contemporáneo}
\textbf{Kaizen} es una filosofía de gestión que propone mejora continua y gradual mediante pequeños cambios incrementales antes que grandes rediseños. Sus principios operativos: (1) \textbf{ciclo PDCA} (\emph{Plan, Do, Check, Act}); (2) \textbf{participación del que ejecuta el proceso} (no decisión solo desde la dirección); (3) \textbf{priorización por matriz impacto/esfuerzo} (foco en lo de alto impacto/bajo esfuerzo primero); (4) \textbf{medición antes y después} (validar que la mejora realmente mejoró). La regla del oficio: \emph{tres mejoras chicas implementadas vencen a diez mejoras grandes propuestas}. Aplicado a procesos digitales (los que mapeaste, automatizaste y mediste con KPIs en este periodo), Kaizen pide cerrar el ciclo: identificar dónde el proceso aún duele, proponer cambios concretos y \emph{empezar mañana}.
\end{softbox}

\begin{softbox}{milcMostaza}{milcGris}{Pregunta-puente}
\textit{¿Cómo sabía la tejedora qué ajuste pequeño de cada manta no era \emph{``perfeccionismo''} sino \emph{aprendizaje real}? ¿Y qué pierde un emprendedor novato cuando, al final del periodo, abandona el proceso \emph{``porque ya funciona''} en lugar de proponer 3 mejoras pequeñas que lo dejen 30\% mejor en un mes?}
\end{softbox}

\begin{softbox}{milcVerde}{milcGris}{Saber y saber hacer}
Al terminar podrás: (1) \textbf{analizar} el proceso del periodo identificando los 3 puntos donde aún duele (cuello de botella residual, error frecuente, fricción del usuario); (2) \textbf{evaluar} cada mejora propuesta en términos de impacto esperado y esfuerzo requerido; (3) \textbf{crear} 3 mejoras concretas con plan de implementación; (4) \textbf{explicar} con tus palabras por qué tres mejoras pequeñas implementadas valen más que diez grandes propuestas.
\end{softbox}

\small
\begin{tabularx}{\textwidth}{>{\bfseries}p{3.0cm}Y}
\rowcolor{milcGradePrimary!12}Autor & Dr. Álvaro Cárdenas Orozco\\
DBA & Diseño y mejora de procesos digitales con criterio ético (MEN, T\&I)\\
Referentes & Pensamiento computacional · Lean / Kaizen · Floridi (ética informacional)\\
Producto final & 1 propuesta de 3 mejoras concretas al proceso del periodo, cada mejora con descripción operativa, impacto esperado (alto/medio/bajo), esfuerzo requerido (alto/medio/bajo), priorización en matriz impacto-esfuerzo y plan de implementación en 2 frases\\
\end{tabularx}
\normalsize

\puente{Antes de empezar, mira el viaje completo. Mapeaste el proceso, lo diagramaste, automatizaste, aplicaste IA, construiste bot y mediste con 5 KPIs. Hoy cierras el ciclo del periodo: \emph{¿qué le falta al proceso para estar mejor el próximo mes?}. La diferencia entre proceso vivo y proceso muerto es Kaizen.}

\begin{titlebox}{milcGradeDark}
Ruta MILC: así vamos a aprender
\end{titlebox}
\begin{tabularx}{\textwidth}{>{\centering\arraybackslash}X>{\centering\arraybackslash}X>{\centering\arraybackslash}X>{\centering\arraybackslash}X}
\phasecard{1}{milcVerde}{milcGris}{Escuta\\[-1mm]\small Observo, escucho y pregunto.} &
\phasecard{2}{milcTurquesa}{milcGris}{Sistematización\\[-1mm]\small Pensamiento computacional.} &
\phasecard{3}{milcMagenta}{milcGris}{Praxis\\[-1mm]\small Construyo y aplico.} &
\phasecard{4}{milcVino}{milcGris}{Evaluación\\[-1mm]\small Reflexión liberadora.}
\end{tabularx}


\puente{Antes de teorizar, vas a usar tu propio proceso del periodo durante 10 minutos como si fueras el usuario final: pasar el Form, recibir el email, conversar con el bot, mirar el dashboard. Y vas a anotar cada fricción real que sientas. Esa lista de fricciones es la materia prima del Kaizen.}

\begin{titlebox}{milcVerde}
Fase 1 · Escuta
\end{titlebox}

\begin{softbox}{milcVerde}{milcGris}{Escena para escuchar}
\textbf{Actividad 1 · ANALIZA --- Usa tu propio proceso} (15 min · individual). Vas a actuar como usuario final de tu proceso del periodo durante 10 minutos: envía el Form, mira si llega bien al Sheet, recibe el email automático, conversa con el chatbot, mira el dashboard. Sin filtrar, anota TODA fricción real que sientas: campo confuso, mensaje feo, espera larga, opción mal ubicada, métrica difícil de leer. Esa lista es tu materia prima.
\end{softbox}

\checkbox Pasé los 10 minutos como usuario, no como diseñador.\\
\checkbox Anoté cada fricción real sin censurarla.\\
\checkbox Identifiqué al menos 5 fricciones concretas en mi propio proceso.

\begin{infoband}{milcVerde}
\textbf{Lo que va al cuaderno (Actividad 1 · ANALIZA).} Título: «Actividad 1 --- Usa tu propio proceso». \textbf{Formato:} tabla 5+ filas × 3 columnas (Fricción detectada~|~Dónde ocurre~|~Impacto en mí como usuario). \textbf{Sabes que terminaste cuando} reconoces honestamente al menos 5 fricciones reales en lo que tú mismo diseñaste.
\end{infoband}

\puente{Ya tienes la lista. Ahora vas a entender la \textbf{anatomía} de una mejora bien propuesta: descripción operativa, impacto esperado, esfuerzo requerido, priorización en matriz, plan en 2 frases. Y los errores típicos: mejoras vagas, mejoras costosas que prometen poco, mejoras sin plan de implementación.}

\begin{titlebox}{milcTurquesa}
Fase 2 · Sistematización con pensamiento computacional
\end{titlebox}

\begin{softbox}{milcTurquesa}{milcGris}{Cuatro pilares aplicados al tema}
Una mejora bien propuesta no parece complicada pero tiene anatomía estricta. Vamos a descomponerla con los pilares del pensamiento computacional para que cada mejora tuya pueda implementarse, medirse y defenderse profesionalmente.
\end{softbox}

\small
\begin{tabularx}{\textwidth}{>{\bfseries\RaggedRight\arraybackslash}p{3.6cm}Y}
\rowcolor{milcTurquesa!90}\color{white}Pilar & \color{white}\bfseries Aplicación al tema\\
1. Descomposición & La mejora se descompone en 5 partes: (1) \textbf{Fricción que resuelve}: 1 frase con el problema real que duele al usuario o al operador del proceso. (2) \textbf{Descripción operativa}: qué cambia concretamente (no \emph{``mejorar UX''}, sino \emph{``mover botón ``escape humano'' al header del chatbot''}). (3) \textbf{Impacto esperado}: alto / medio / bajo, justificado con número o regla. (4) \textbf{Esfuerzo requerido}: alto / medio / bajo, en horas-persona. (5) \textbf{Plan en 2 frases}: \emph{``Esta semana: X. Próxima semana: Y''}.\\
2. Reconocimiento de patrones & Toda priorización sigue la \textbf{matriz impacto/esfuerzo} (2×2): cuadrante \emph{alto impacto / bajo esfuerzo} se hace primero (\emph{quick wins}); \emph{alto impacto / alto esfuerzo} se planifica para periodo siguiente; \emph{bajo impacto / bajo esfuerzo} se hace si sobra tiempo; \emph{bajo impacto / alto esfuerzo} se descarta. Esta matriz es la herramienta Kaizen más simple y poderosa.\\
3. Abstracción & Lo esencial cabe en una regla: \emph{tres mejoras implementadas vencen a diez mejoras propuestas}. La medida del Kaizen no es la cantidad de ideas; es la cantidad de \emph{ciclos PDCA cerrados}. Cada ciclo: planificar la mejora, hacerla, verificar con KPI si funcionó, ajustar. Sin medición antes-después, no hay Kaizen real, solo cambios.\\
4. Algoritmo conceptual & Algoritmo: (1) ordena las 5+ fricciones de Actividad 1 por gravedad; (2) toma las 3 más graves; (3) redacta una mejora por cada una con las 5 partes; (4) ubica cada mejora en la matriz impacto-esfuerzo; (5) elige el orden de implementación (quick wins primero); (6) escribe el plan en 2 frases; (7) define qué KPI medirá el éxito de cada mejora.\\
\end{tabularx}
\normalsize

\begin{drawbox}{milcTurquesa}{Anatomía de una mejora bien propuesta}
\textbf{Fricción:} qué duele hoy. \\[1mm] \textbf{Cambio operativo:} qué se hace concretamente. \\[1mm] \textbf{Impacto esperado:} alto/medio/bajo + justificación. \\[1mm] \textbf{Esfuerzo:} alto/medio/bajo + horas-persona. \\[1mm] \textbf{Plan:} \emph{esta semana / próxima semana}. \\[1mm] \textbf{KPI que medirá el éxito:} cuál de los 5 KPIs subirá o bajará. \\[1mm] \textbf{Cuadrante:} quick win / planificar / si sobra / descartar.
\end{drawbox}

\begin{softbox}{milcMostaza}{milcGris}{Errores comunes y su corrección}
(1) Mejora vaga (\emph{``mejorar comunicación''}): no se puede implementar ni medir. (2) Mejora cara que promete poco: cuadrante \emph{bajo impacto / alto esfuerzo}, debería descartarse. (3) Mejora sin plan: queda como idea bonita en la libreta. (4) Mejora sin KPI asociado: no se puede verificar si funcionó. (5) Proponer 15 mejoras: la energía se dispersa; mejor 3 bien ejecutadas.
\end{softbox}

\begin{infoband}{milcTurquesa}
\textbf{Lo que va al cuaderno (Actividad 2 · EVALÚA).} Título: «Actividad 2 --- Anatomía de la mejora». \textbf{Formato:} ficha con las 5 partes + matriz 2×2 dibujada con ejes \emph{impacto} (vertical) y \emph{esfuerzo} (horizontal) + 5 errores comunes con marca 🚨. \textbf{Sabes que terminaste cuando} podrías clasificar mejoras ajenas en el cuadrante correcto.
\end{infoband}

\puente{Tienes el método. Toca aplicarlo: priorizar 3 de las fricciones identificadas, redactar la mejora correspondiente, situarla en la matriz impacto-esfuerzo y escribir el plan en 2 frases.}

\begin{titlebox}{milcMagenta}
Fase 3 · Praxis con pensamiento computacional
\end{titlebox}

\begin{softbox}{milcMagenta}{milcGris}{Construyendo el producto con los cuatro pilares}
\textbf{Actividad 3 · CREA + EXPLICA --- 3 mejoras priorizadas} (30 min · individual). Hora de crear. Vas a tomar las 5+ fricciones de Actividad 1, elegir las 3 más graves, redactar la mejora correspondiente con las 5 partes, ubicarla en la matriz y escribir el plan de implementación.
\end{softbox}

\small
\begin{tabularx}{\textwidth}{>{\bfseries\RaggedRight\arraybackslash}p{3.6cm}Y}
\rowcolor{milcMagenta!90}\color{white}Pilar & \color{white}\bfseries Aplicación al producto\\
Descomposición & Tu producto se descompone en 4 piezas: (1) lista priorizada de las 3 mejoras elegidas; (2) ficha completa por mejora con las 7 partes de anatomía (fricción, cambio operativo, impacto, esfuerzo, plan, KPI, cuadrante); (3) matriz impacto-esfuerzo dibujada con las 3 mejoras ubicadas; (4) orden de implementación con justificación.\\
Patrones & Sigue el patrón \textbf{``quick wins primero''}: empieza por lo de alto impacto / bajo esfuerzo, no por lo más ambicioso. El Kaizen valora el \emph{movimiento real} sobre la \emph{idea grande}. Tres quick wins en un mes producen mejor proceso que un rediseño que nunca se ejecuta.\\
Abstracción & Cada plan se redacta en 2 frases concretas: \emph{``Esta semana: rediseño el menú del chatbot reordenando los 3 botones más usados arriba. Próxima semana: pruebo con 5 usuarios y mido tasa de finalización contra el dashboard.''}. Esa concreción separa la idea del compromiso.\\
Algoritmo de armado & Algoritmo: (1) prioriza tus fricciones; (2) toma las 3 más graves; (3) redacta cada mejora con las 7 partes; (4) ubícalas en la matriz; (5) ordena por cuadrante (quick wins → planificar); (6) escribe el plan; (7) define el KPI que medirá el éxito; (8) comparte con un compañero para que critique las propuestas.\\
\end{tabularx}
\normalsize

\begin{drawbox}{milcMagenta}{Checklist antes de entregar las 3 mejoras}
\checkbox 5+ fricciones reales identificadas en Actividad 1. \\[1mm] \checkbox 3 mejoras priorizadas con las 7 partes de anatomía. \\[1mm] \checkbox Matriz impacto-esfuerzo dibujada con las 3 ubicadas. \\[1mm] \checkbox Orden de implementación con justificación (quick wins primero). \\[1mm] \checkbox Plan en 2 frases por mejora. \\[1mm] \checkbox KPI asociado a cada mejora para validar el ciclo PDCA.
\end{drawbox}

\begin{softbox}{milcMagenta}{milcGris}{Plantilla de trabajo}
\textbf{Proceso mejorado.} \textit{Nombre + estado actual + 5 KPIs base.} \\[1mm] \textbf{Mejora 1 (quick win).} \textit{Las 7 partes.} \\[1mm] \textbf{Mejora 2.} \textit{Lo mismo.} \\[1mm] \textbf{Mejora 3.} \textit{Lo mismo.} \\[1mm] \textbf{Matriz.} \textit{Dibujo con las 3 ubicadas.} \\[1mm] \textbf{Plan global.} \textit{Mes 1: mejora A. Mes 2: B y C.}
\end{softbox}

\begin{infoband}{milcMagenta}
\textbf{Lo que va al cuaderno (Actividad 3 · CREA + EXPLICA).} Título: «Actividad 3 --- 3 mejoras priorizadas». \textbf{Formato:} 3 fichas (una por mejora) + matriz dibujada + plan global. \textbf{Extensión:} 1-2 páginas de cuaderno. \textbf{Sabes que terminaste cuando} compartiste las 3 mejoras con un compañero o emprendedor y él podría empezar a ejecutar la mejora 1 esta misma semana sin pedir explicación extra.
\end{infoband}

\puente{Lo que sigue resume \emph{qué} entregas hoy: 3 mejoras priorizadas + matriz impacto-esfuerzo dibujada + plan de implementación por mejora. El criterio principal: que un emprendedor que ejecute las 3 mejoras en 1 mes pueda \textbf{medir la diferencia} con los KPIs que ya tienes.}

\begin{titlebox}{milcMostaza}
Producto de la sesión
\end{titlebox}
\textbf{3 mejoras priorizadas con matriz impacto-esfuerzo y plan}.
\begin{softbox}{milcMostaza}{milcGris}{Criterios de calidad}
(1) 5+ fricciones identificadas a partir de uso real. (2) 3 mejoras con las 7 partes de anatomía. (3) Matriz impacto-esfuerzo con ubicación justificada. (4) Quick wins identificados y priorizados. (5) Plan de implementación en 2 frases por mejora.
\end{softbox}

\puente{Antes de cerrar, mira las mejoras desde las cinco dimensiones humanas. Mejorar bien es respeto por el oficio que sostiene el proceso; mejorar mal es agregar trabajo a quienes ya cargan con lo invisible.}

\begin{titlebox}{milcVino}
Fase 4 · Evaluación liberadora — cinco dimensiones
\end{titlebox}

\begin{softbox}{milcVino}{milcGris}{Sentido de la evaluación}
Evaluar no es solo revisar una nota. Es reconocer qué aprendí, qué cambió en mí, qué emociones manejé, cómo me relaciono con la ciudad y la comunidad, y qué le dejaré a quien viene detrás.
\end{softbox}

\textbf{Mi reflexión liberadora en cinco dimensiones.} Completa cada frase iniciada:

\small
\begin{tabularx}{\textwidth}{>{\bfseries\RaggedRight\arraybackslash}p{3.4cm}Y>{\RaggedRight\arraybackslash}p{4.6cm}}
\rowcolor{milcVino}\color{white}Dimensión & \color{white}\bfseries Mi respuesta (frame starter) & \color{white}\bfseries Algo que descubrí\\
1. Personal & Lo que más cambió en mí fue \linead{6.5cm} & \linead{4.2cm}\\
2. Emocional & La emoción más fuerte fue \linead{2.5cm} y la regulé \linead{3cm} & \linead{4.2cm}\\
3. Ciudadana & En democracia esto importa porque \linead{6cm} & \linead{4.2cm}\\
4. Local (Cartago) & En mi barrio sería útil para \linead{6.5cm} & \linead{4.2cm}\\
5. Intergeneracional & A mi abuela/abuelo le diría \linead{6.5cm} & \linead{4.2cm}\\
\end{tabularx}
\normalsize

\begin{infoband}{milcVino}
\textbf{Para la próxima sustentación.} Vuelve a esta tabla en seis meses. Verás que las respuestas cambian — no porque lo aprendido cambie, sino porque tú cambias. La evaluación liberadora es una conversación contigo a través del tiempo.
\end{infoband}


\puente{Y para terminar, tres voces sobre Kaizen. Dussel sobre las mejoras que solo benefician al dueño, Marco Aurelio sobre la paciencia del oficio que se afina, Floridi sobre la mejora continua como ética profesional contemporánea.}

\begin{titlebox}{milcGradeDark}
Triángulo de pensamiento · cierre filosófico
\end{titlebox}

\begin{softbox}{milcVino}{milcGris}{Dussel · lente del nosotros}
\textit{``Toda mejora decide a quién se le ahorra trabajo y a quién se le agrega.''}

Una mejora que reduce el tiempo del cliente puede aumentar la carga del operador (más mensajes, más trabajo invisible). Una automatización que ahorra al dueño puede dejar sin trabajo a quien hacía esa tarea. La phronesis del Kaizen exige preguntar siempre: \emph{¿quién paga el costo de esta mejora?, ¿a quién se le agrega trabajo invisible?}. Las mejoras éticas distribuyen los beneficios; las mejoras egoístas concentran ganancias y descargan costos.

\textbf{En mis 3 mejoras, ¿quién gana tiempo y quién paga el costo escondido? ¿Hay alguna mejora donde el costo cae sobre un trabajador o usuario que ya estaba sobrecargado?}
\end{softbox}
\linead{\linewidth}\\[1mm]
\linead{\linewidth}

\begin{softbox}{milcMostaza}{milcGris}{Estoicismo · Marco Aurelio · cuidado interior}
\textit{``El oficio nunca se termina de aprender; se afina pieza tras pieza.''}

Es tentador buscar el rediseño grande que arregle todo de un golpe. La sabiduría estoica (que coincide con Kaizen) recomienda lo contrario: \emph{tres mejoras chicas implementadas valen más que diez grandes propuestas}. La virtud del oficio está en el ciclo PDCA repetido con paciencia, no en el rediseño espectacular. Esa paciencia se entrena al cerrar 3 ciclos al mes y no perseguir el cambio absoluto.

\textbf{¿Estoy diseñando mejoras pequeñas implementables o estoy proponiendo un rediseño grande que nunca se ejecutará? ¿Cuántos ciclos PDCA cerré este periodo?}
\end{softbox}
\linead{\linewidth}\\[1mm]
\linead{\linewidth}

\begin{softbox}{milcTurquesa}{milcGris}{Floridi · lente de la infoesfera}
\textit{``La mejora continua es la ética profesional del oficio digital contemporáneo.''}

Toda organización contemporánea que merezca seguir existiendo aplica algún ciclo de mejora continua. No es opcional: es la ética profesional mínima del XXI. Aprender Kaizen a los 16 años es alfabetización de oficio: aceptar que el producto nunca está terminado, que la responsabilidad nunca se acaba, que la mejora es continua porque la realidad cambia continuamente.

\textbf{¿Asumo mi proceso del periodo como \emph{``ya terminé''} o como \emph{``ahora empieza la mejora continua''}? ¿Qué cambia en mi compromiso si acepto que el oficio nunca termina?}
\end{softbox}
\linead{\linewidth}\\[1mm]
\linead{\linewidth}

\begin{titlebox}{milcVino}
Compromiso final
\end{titlebox}

\small
\begin{tabularx}{\textwidth}{>{\RaggedRight\arraybackslash}p{3.0cm}YYY}
\rowcolor{milcVino}\color{white}\bfseries Criterio & \color{white}\bfseries Lo logré & \color{white}\bfseries En proceso & \color{white}\bfseries Necesito apoyo\\
Comprensión del tema & Explico ideas centrales con mis palabras. & Reconozco ideas pero falta precisión. & Necesito repasar conceptos base.\\
Pensamiento computacional & Apliqué los 4 pilares al diseñar mi producto. & Apliqué algunos pilares. & Necesito releer la fase 2.\\
Aplicación y producto & Mi producto cumple los criterios y se puede revisar. & Producto funcional con ajustes pendientes. & Aún en construcción.\\
Reflexión 5 dimensiones & Respondí cada dimensión con honestidad. & Respondí algunas con detalle. & Mis respuestas son muy generales.\\
Triángulo de pensamiento & Conecté las 3 voces con mi trabajo real. & Conecté al menos una. & No relaciono las voces con mi proceso.\\
\end{tabularx}
\normalsize

\textbf{Hoy entendí que:}\\[1mm]
\linead{\linewidth}\\[1mm]
\linead{\linewidth}

\textbf{Mi compromiso para la próxima sesión es:}\\[1mm]
\linead{\linewidth}\\[1mm]
\linead{\linewidth}

\begin{softbox}{milcMostaza}{milcGris}{Cita de despedida · Marco Aurelio}
\textit{``El obstáculo en el camino se vuelve el camino. Lo que estorba al camino se vuelve camino.''}\\[1mm]
Cada error de hoy, bien observado, se convierte en pieza del camino que pisarás mañana.
\end{softbox}

\small
\begin{tabularx}{\textwidth}{>{\bfseries}p{3.6cm}Y}
\rowcolor{milcGradeDark!12}Nombre completo: & \linead{10cm}\\
Fecha: & \linead{4cm} \quad Curso/grupo: \linead{4cm}\\
Mi firma: & \linead{9cm}\\
\end{tabularx}
\normalsize

\begin{infoband}{milcGradeDark}
\textbf{Para la próxima semana.} Implementa la mejora 1 (quick win) y mide con el KPI asociado antes y después. Si el KPI cambió favorablemente, documenta el caso; si no, ajusta la mejora o reemplázala. La phronesis del Kaizen se entrena cerrando ciclos PDCA, no proponiendo planes que nunca se ejecutan.
\end{infoband}

\end{document}
